\uuid{RB17}
\exo7id{3813}
\titre{exo7 3813}
\auteur{quercia}
\organisation{exo7}
\datecreate{2010-03-11}
\isIndication{false}
\isCorrection{true}
\chapitre{Espace euclidien, espace normé}
\sousChapitre{Problèmes matriciels}
\cpge{mpsi-mp2i.calcul-matriciel-systemes.e, pcsi.determinants.c, ptsi.determinants.a, mpsi-mp2i.espaces-prehilbertiens.b, pcsi.espaces-prehilbertiens.b, ptsi.geometrie-espace.b, bcpst1.matrices, bcpst1.geometrie.f, mp-mpi.endomorphismes-euclidien.b, pc.endomorphismes-euclidiens.b, psi.endomorphismes-euclidien.b, pt.isometries-euclidien.b, mp-mpi.endomorphismes-euclidien.c, pc.endomorphismes-euclidiens.a, psi.endomorphismes-euclidien.a, pt.isometries-euclidien.a, bcpst2.produit-scalaire.a}
\module{Algèbre}
\niveau{L2}
\difficulte{}

\contenu{
\question{
Déterminer $a$, $b$, $c$, réels non nuls pour que la matrice
$M = -\frac23\begin{pmatrix}-\frac12&\frac ac &\frac ab\\[3pt]
                       \frac ca &-\frac12 &\frac ca\\[3pt]
                       \frac ba &\frac ac &-\frac12\end{pmatrix}$
soit la matrice d'une isométrie.
}
\reponse{
$a=b=\pm c$.
}
}
